\par\noindent\textbf{(1)} 本地竖直方向的地球自转分量为 \(\Omega_z=\Omega_E\sin\lambda\)。小摆的自由固有频率 \(\omega_0=\sqrt{g/l}\)，水平动力学为 \(\ddot x+\omega_0^2x=2\Omega_z\dot y\)、\(\ddot y+\omega_0^2y=-2\Omega_z\dot x\)，舍弃 \(O(\Omega_E^2)\) 离心修正。
\par\noindent\textbf{(2)} 定义 \(z=x+iy\)，两方程合并 \(\ddot z+2i\Omega_z\dot z+\omega_0^2z=0\)。试 \(z\propto e^{-i\omega t}\) 得 \(\omega^2-2\Omega_z\omega-\omega_0^2=0\)，根为 \(\omega=\Omega_z\pm\sqrt{\omega_0^2+\Omega_z^2}\)，分别对应两种相反圆振动。
\par\noindent\textbf{(3)} 令 \(z=e^{-i\Omega_z t}u\)，则 \(\ddot u+(\omega_0^2+\Omega_z^2)u=0\)。因此在近似惯性参考的 \(u\) 中摆动平面固定，在地球共转坐标中以 \(\dot\psi=-\Omega_E\sin\lambda\) 慢旋；北半球为顺时针进动。
\par\noindent\textbf{(4)} 极点 \(|\sin\lambda|=1\) 时进动周期为恒星日；赤道 \(\sin\lambda=0\) 时一阶Coriolis慢进动消失。一般纬度周期 \(T_p=2\pi/(|\Omega_E\sin\lambda|)\)，真实装置还受空气阻尼、悬挂非对称与驱动影响。
